\documentclass[11pt]{article}
\usepackage[T1]{fontenc}
\usepackage{lmodern,amsmath,amssymb,amsthm,mathtools}
\usepackage[margin=1in]{geometry}
\usepackage{microtype}
\usepackage[hidelinks]{hyperref}
\hypersetup{pdftitle={A positive-density family of dihedral fields with Bloch lattice index five},pdfauthor={},pdfcreator={},pdfproducer={}}
\pdfinfoomitdate=1
\pdftrailerid{}
\pdfsuppressptexinfo=15
\newtheorem{theorem}{Theorem}
\newtheorem{lemma}[theorem]{Lemma}
\newtheorem{proposition}[theorem]{Proposition}
\theoremstyle{remark}\newtheorem{remark}[theorem]{Remark}
\newcommand{\Q}{\mathbb Q}\newcommand{\Z}{\mathbb Z}
\newcommand{\F}{\mathbb F}\newcommand{\Gal}{\operatorname{Gal}}
\newcommand{\res}{\operatorname{res}}\newcommand{\cor}{\operatorname{cor}}
\newcommand{\et}{\mathrm{\acute et}}
\newcommand{\coker}{\operatorname{coker}}
\newcommand{\N}{\operatorname{N}}
\newcommand{\dgen}{\operatorname{d}}
\title{A positive-density family of dihedral fields\protect\\with Bloch lattice index five}
\author{}\date{}
\begin{document}
\maketitle
\begin{abstract}
We construct an infinite family of non-real dihedral extensions of degree ten over
\(\Q\), with fixed quadratic subfield \(\Q(i)\), for which the Bloch lattice index
\(Q_2\) is five. A fixed Chebotarev set of rational primes has density \(1/50\).
For each parameter prime in this set, at least two of the four mixed cyclic
quintic ray class extensions of conductor \((25\ell)\) have the asserted index.
Their five-primary tame kernels are \((\Z/5)^2\). The proof combines a
two-dimensional cup-product pencil, the first generalized Bockstein map, integral
co-descent, and a motivic genus formula. It does not assume a Lichtenbaum special
value formula or determine a Bloch lattice by numerical regulators.
\end{abstract}

\section{The index and the statement}
For a number field \(J\), let \(\mathcal B(J)\) denote the free quotient of its
Bloch group, equivalently the free quotient of \(K_3^{\mathrm{ind}}(J)\).
These identifications are natural in field embeddings; the regulator identifies
the free quotient with its Bloch lattice. Let \(M/\Q\) be a non-real
\(D_{10}\)-extension, let \(k\) be its quadratic subfield, and let
\(F_1,F_2\) be two different quintic reflection subfields. The index considered
by Browkin--Gangl and Guo--Qin is
\[
 Q_2(M)=[\mathcal B(M):\mathcal B(k)+\mathcal B(F_1)+\mathcal B(F_2)],
\]
where subgroups are understood via their natural images. Guo--Qin
\cite[Theorem 4.3]{GQ} prove \(Q_2(M)\mid5\).

Put \(k=\Q(i)\), \(E=k(\mu_5)=\Q(\zeta_{20})\), and let \(S_0\) consist of
the primes above five and the infinite places. All sets of primes are pulled
back to every extension in use. Write \(K_2(O_J)(5)\) for the five-primary
subgroup. The base vanishing needed below has a short verification.

\begin{lemma}\label{lem:base}
\(K_2(\Z[i])(5)=0\), and \(H^2(G_{k,S_0},\F_5(2))=0\).
\end{lemma}
\begin{proof}
The field \(E=\Q(\zeta_{20})\) has class number one; the accompanying
fixed-field certificate verifies this fact independently. Kummer theory gives
\(H^2(O_E[1/5],\mu_5)=\operatorname{Br}(O_E[1/5])[5]\).
The Brauer invariant sequence identifies this group with the kernel of
\(\F_5^2\xrightarrow{+}\F_5\), one coordinate for each five-prime.
Both primes are fixed by \(\Delta=\Gal(E/k)\), and local invariants show
that \(\Delta\) acts trivially on this kernel. Twisting by
\(\F_5(1)\) therefore gives a nontrivial cyclotomic character on
\(H^2(G_{E,S_0},\F_5(2))\). Its \(\Delta\)-invariants vanish.
Since \(|\Delta|=4\), Hochschild--Serre identifies those invariants with
\(H^2(G_{k,S_0},\F_5(2))\). The integral weight-two comparison with the
tame kernel, reduction modulo five, and finiteness then give the asserted
five-primary vanishing.
\end{proof}

\begin{theorem}\label{thm:main}
There is a set \(\mathcal P\) of rational primes of natural and Dirichlet density
\(1/50\), all congruent to one modulo twenty, with the following property.
For each \(\ell\in\mathcal P\), at least two of the four anti-conjugate mixed
cyclic quintic ray class extensions \(M/k\) of conductor \((25\ell)\) satisfy
\[
 \Gal(M/\Q)\simeq D_{10},\qquad Q_2(M)=5,
 \qquad K_2(O_M)(5)\simeq(\Z/5)^2.
\]
Their discriminants are
\(\operatorname{disc}(M)=-2^{10}5^{16}\ell^8\).
If \(a=v_5(\ell-1)\) and \(S=S_0\cup\{v:v\mid\ell\}\), then
\[
 H^2(G_{M,S},\Z_5(2))\simeq(\Z/5^{a+1})^2.
\]
In particular, there are infinitely many pairwise non-isomorphic such fields
with the same quadratic subfield and the same degree over \(\Q\).
\end{theorem}

The theorem permits a choice among four actual ray class fields for each
parameter. It does not assert that an arbitrary preselected ray has index five.
The density is a density of parameter primes, not a density ordered by field
discriminant.

\section{The fixed extensions and their Chebotarev set}
The five-primary character space of the ray class group of \(k\) modulo
\((25)\) has dimension two. Complex conjugation exchanges the two primes above
five. Its negative eigenspace has dimension one. Fix a nonzero character
\(\eta_0\) in this eigenspace, with fixed field \(M_0\). Thus \(M_0/k\) is
cyclic of degree five and conductor \((25)\), and \(M_0/\Q\) is dihedral.
This uses \(h(k)=1\) and \(|\mu(k)|=4\).

Let \(T=\F_5(2)\). The integral group
\(H^1(G_{k,S_0},\Z_5(2))\) is free of rank one over \(\Z_5\): its rank is
the Borel rank \(r_2(k)=1\), and its torsion is
\(H^0(k,\Q_5/\Z_5(2))=0\), since \(\sqrt5\notin k\).
The vanishing of the integral \(H^2\) and Bockstein give
\[
 H^1(G_{k,S_0},\Z_5(2))/5\simeq H^1(G_{k,S_0},T).
\]
Restriction to \(E\) is injective, because \([E:k]=4\) is prime to five.
The Kummer comparison
\(H^1(E,T)=H^1(E,\mu_5)\otimes\F_5(1)\) turns its image into a
one-dimensional nonzero subgroup \(D\subset E^\times/E^{\times5}\).
This is the weight-two Tate line of \cite{AM19}. It is unramified outside five
in the sense that \(E_T=E(\sqrt[5]{D})/E\) is unramified outside five.

Let \(\Delta=\Gal(E/k)\simeq C_4\), and let \(\omega\) be its cyclotomic
character modulo five. The line \(D\) has character \(\omega^{-1}\).
Consequently, the Kummer pairing with target \(\mu_5\) shows that
\(\Delta\) acts on \(\Gal(E_T/E)\) with character \(\omega^2\).
On \(\Gal(M_0E/E)\), the action is trivial, since \(M_0E/k\) is abelian.
Thus the two cyclic quintic extensions of \(E\) are different and linearly
disjoint. The Tate line is also stable under conjugation over \(\Q\);
both factors, and hence
\[
 L=E_T M_0,
\]
are normal over \(\Q\). We have \([L:\Q]=8\cdot25=200\).

Let \(\mathcal P\) consist of the unramified rational primes whose Frobenius
in \(L/\Q\) is trivial on \(E_T\) and nontrivial on \(M_0E/E\).
In \(\Gal(L/E)=C_5\times C_5\), this condition selects exactly four elements.
It is invariant under conjugation, because both factors are normal over
\(\Q\). Chebotarev gives density \(4/200=1/50\).
Such primes split completely in \(E\), so \(\ell\equiv1\pmod{20}\).
For either prime of \(k\) above \(\ell\), the unramified character
\(\eta_0\) has nonzero Frobenius value, and the Tate Frobenius is zero.
Different choices of a prime of \(E\) change Frobenius by a nonzero scalar
conjugation; its vanishing and the dimension of its span are unchanged.

\section{A cup-product pencil on the fixed base field}
Fix \(\ell\in\mathcal P\), and let \(v_1,v_2\) be the primes of \(k\) over
\(\ell\). The anti-conjugate five-order character space modulo \((\ell)\)
is one-dimensional. Choose a nonzero character \(\eta_\ell\) there.
For \(\lambda\in\F_5^\times\), put
\[
 \eta_\lambda=\eta_0+\lambda\eta_\ell,
 \qquad M_\lambda=\overline{k}^{\ker\eta_\lambda}.
\]
The two characters have different ramification sets and are independent.
These are four different cyclic quintic layers, each fully ramified at both
five-primes and both \(\ell\)-primes and with conductor \((25\ell)\).
Conjugation negates each character, so each field is dihedral over \(\Q\).
The conductor--discriminant formula gives
\[
 \operatorname{disc}(M_\lambda)
 =(-4)^5\N_{k/\Q}((25\ell)^4)
 =-2^{10}5^{16}\ell^8.
\]

Set \(\mathcal G=G_{k,S}\), with \(S=S_0\cup\{v_1,v_2\}\).
The base group \(H^2(G_{k,S_0},T)\) is zero. Localization and purity give
\begin{align}
 H^1(\mathcal G,T)&\longrightarrow
 \bigoplus_{j=1}^2 H^0(\kappa(v_j),\F_5(1))\longrightarrow0,\label{eq:inertia}\\
 H^2(\mathcal G,T)&\xrightarrow{\sim}
 \bigoplus_{j=1}^2 H^2(k_{v_j},T)\simeq\F_5^2.\label{eq:h2loc}
\end{align}
In \eqref{eq:inertia}, each summand is \(\F_5\), since
\(\ell\equiv1\pmod5\), and the map records tame inertia residues.
Surjectivity follows from the preceding zero \(H^2(G_{k,S_0},T)\).
For \eqref{eq:h2loc}, localization identifies the new finite-field
\(H^1(\kappa(v_j),\F_5(1))\) with the local \(H^2\); the next global
\(H^3\) is zero. The five-adic local \(H^2\) groups vanish because those
completions do not contain \(\mu_5\).

Choose \(y_1,y_2\in H^1(\mathcal G,T)\) mapping to the standard basis in
\eqref{eq:inertia}. The character \(\eta_0\) is unramified at these primes,
with Frobenius values \(u,-u\), where \(u\ne0\). Locally, cup product pairs
an unramified \(\F_5\)-character perfectly with the one-dimensional tame
inertia residue in \(H^1(T)\); the cup of two unramified classes is zero.
Consequently \(\eta_0\cup-\), restricted to the span of \(y_1,y_2\), has
an invertible diagonal matrix \(C_0\) in \eqref{eq:h2loc}.

Let \(C_\ell\) be the matrix for \(\eta_\ell\cup-\) on the same two
vectors. For \(\eta_\lambda\), the matrix is
\[
 C_\lambda=C_0+\lambda C_\ell.
\]
Its determinant is a polynomial of degree at most two over \(\F_5\), with
nonzero constant coefficient. It has at most two roots. Therefore at least
two of the four nonzero \(\lambda\) satisfy
\begin{equation}\label{eq:cup}
 \eta_\lambda\cup-:H^1(\mathcal G,T)\twoheadrightarrow H^2(\mathcal G,T).
\end{equation}
No assertion about the distribution of residue symbols of varying prime
generators is required: \(C_\ell\) may depend arbitrarily on \(\ell\).

\section{From cup product to integral co-descent}
Fix a line satisfying \eqref{eq:cup}. Write \(M=M_\lambda\),
\(\mathcal N=\ker\eta_\lambda=G_{M,S}\), and
\(G=\mathcal G/\mathcal N=\langle\rho\rangle=C_5\).
Put \(\Omega=\F_5[G]\), \(I=(\rho-1)\), and \(A=H^2(\mathcal N,T)\).
Theorem 2.2.4 of Lam--Liu--Sharifi--Wake--Wang \cite{LLSWW}, with
\(R=\F_5\), \(n=1\), and finite \(H=G\), gives
\[
 IA/I^2A\simeq\coker\bigl(H^1(\mathcal G,T)
 \xrightarrow{\,\eta_\lambda\cup-\,}H^2(\mathcal G,T)\bigr).
\]
All required flatness conditions hold over \(\F_5\).
For finite \(H\), its Iwasawa cohomology is ordinary subgroup cohomology
\cite[Remark 2.2.1]{LLSWW}; the first Bockstein is the cup product
\cite[Proposition 2.3.3]{LLSWW}.
Equation \eqref{eq:cup} gives \(IA=I^2A\). Since \(I^5=0\), it gives
\(IA=0\). The natural top-degree co-descent isomorphism
\(A_G\simeq H^2(\mathcal G,T)\) therefore proves \(A\simeq\F_5^2\).

Put
\[
 B=H^2(G_{M,S},\Z_5(2)),\quad W=K_2(O_M)(5),
 \quad a=v_5(\ell-1),\quad R=(\Z/5^a)^2.
\]
These integral cohomology groups are finite. Cohomological dimension at most
two and the Mittag--Leffler property for finite coefficient groups give
continuous \(H^3(G_{M,S},\Z_5(2))=0\). Thus Bockstein gives
\(B/5B\simeq A\), and \(\dgen(B)=2\).

The integral localization sequence is
\begin{equation}\label{eq:integralloc}
 0\longrightarrow W\longrightarrow B
 \xrightarrow{\lambda_{\mathrm{loc}}}R\longrightarrow0.
\end{equation}
The base group \(H^2(G_{k,S},\Z_5(2))\) is also \(R\), because its tame
kernel has zero five-primary part. Integral co-descent
\cite[Theorem 1.5 and the following paragraph]{AM12} identifies
\(\cor:B_G\xrightarrow{\sim}R\).
The localization diagram is natural: on residue multiplicative groups the
corestriction is the residue-field norm, while restriction is the
ramification-index power. Here \(f=1\) and \(e=5\), so, in matched residue
coordinates,
\begin{equation}\label{eq:transfer}
 q:=\cor=\lambda_{\mathrm{loc}},\qquad
 (\rho-1)B=W,\qquad \res\,q=N_G.
\end{equation}
The first equality uses the residue norm, not multiplication by the
ramification index.

\begin{lemma}\label{lem:integral}
With the preceding notation, \(W=5^aB\), \(5^{a+1}B=0\),
\(N_G=5\operatorname{id}_B\), and \(G\) acts trivially on \(W\).
\end{lemma}
\begin{proof}
A surjection from a finite abelian five-group with two generators onto
\((\Z/5^a)^2\) forces both cyclic factors to have exponent at least \(a\).
Hence \(|B/5^aB|=|R|\). Since \(5^aB\subset\ker q\), comparison of orders
in \eqref{eq:integralloc} gives \(W=5^aB\).
Let \(t=\rho-1\). Then \(tB=5^aB\) and \(t^jB=5^{aj}B\).
The identity
\[
 0=5t+10t^2+10t^3+5t^4+t^5
\]
implies \(5^{a+1}B\subset5^{a+2}B\), as every other term has image contained
in the latter group when \(a\ge1\). Finiteness forces \(5^{a+1}B=0\).
Thus \(t^2B=0\); in
\(N_G=5+10t+10t^2+5t^3+t^4\), all nonconstant terms vanish, including
\(10tB\subset5^{a+1}B\). Finally
\(tW=5^atB=5^{2a}B=0\).
\end{proof}

\section{The genus formula and the actual Bloch lattice}
The motivic tame-kernel genus formula
\cite[Proposition 4.11 and Corollary 4.12]{AM19} gives
\[
 \frac{|W_G|}{|K_2(O_k)(5)|}=5^{2-t_2}.
\]
Here \(t_2\) is the dimension of the span of the two tame Frobenius elements
in the fixed extension \(E_T/E\). It depends only on the weight-two Tate line
of \(k\) and the two tame ramified primes. All four mixed lines have the same
tame set and nonzero local characters; changing \(\lambda\) multiplies a tame
character by a nonzero scalar and does not change its norm kernel.
By construction of \(\mathcal P\), \(t_2=0\) for all four lines.

Lemma \ref{lem:integral} makes \(W_G=W\). The genus formula gives \(|W|=25\),
and \(5W=0\). Both cyclic factors of \(B\) must consequently have exponent
\(a+1\). The equality \(\res q=5\operatorname{id}_B\), with \(q\) surjective,
gives \(\res(R)=5B\). These groups have equal order, so restriction is
injective. For a finite cyclic action \(|B^G|=|B_G|=|R|\), so restriction is
also surjective onto \(B^G\). In particular the integral restriction
\begin{equation}\label{eq:f2}
 f_2:H^2(G_{k,S},\Z_5(2))\longrightarrow B^G
\end{equation}
has zero cokernel.

We include the lattice consequence to make the link to \(Q_2\) explicit.
It uses no special-value conjecture.
\begin{proposition}\label{prop:lattice}
For a non-real \(D_{10}\)-extension \(M/\Q\), if \eqref{eq:f2} has zero
cokernel for a set \(S\) covering its relative ramification and all five-primes,
then \(Q_2(M)=5\).
\end{proposition}
\begin{proof}
Let \(\Lambda=K_3^{\mathrm{ind}}(M)\otimes\Z_5\).
It is a free lattice of rank five. The torsion formula for indecomposable
\(K_3\) and the Chern comparison \cite[Theorem 4.5 and Corollary 4.6]{Levine}
show that five-torsion would require \(\Q(\sqrt5)\subset M\).
But the unique quadratic subfield of a non-real dihedral field is imaginary.
Thus there is no five-torsion.
Levine's integral Galois descent \cite[Theorem 4.13]{Levine} identifies the
actual subfield images with \(\Lambda^{C_5}\) and the two reflection invariant
lattices. Its finite-quotient order condition is prime to the field
characteristic, which is zero; it does not exclude degree five.
The natural Suslin exact sequence identifies the free \(K_3\) quotient with
the free Bloch quotient, compatibly with these images.

Weight-two integral \(H^1\) does not increase on enlarging \(S\), because
the possible new finite-field \(H^0(\kappa(v),\Z_5(1))\) is zero.
The comparison and \cite[Theorem 1.5]{AM12} give
\[
 \coker f_2=H^2(C_5,\Lambda)
 =\Lambda^{C_5}/N_G\Lambda=0.
\]
Write \(\Lambda^{C_5}=\Z_5y\), a saturated rank-one sublattice. Its quotient
is a rank-one free module over \(\Z_5[\zeta_5]\). Lift its generator to \(x\),
and write \(N_Gx=by\). The zero norm quotient forces \(b\) to be a unit;
thus \(\Lambda\simeq\Z_5[C_5]x\).

Let \(c\) be a reflection. On the invariant line it acts as \(-1\), since
that line descends to the imaginary quadratic field, whose \(K_3\) free
quotient has no rational invariant part. In the regular group ring
\(\mathcal R=\Z_5[C_5]\), write \(cx=wx\), so \(w\bar w=1\) and
\(\operatorname{aug}(w)=-1\), with \(\bar\rho=\rho^{-1}\).
The unit \(u=(1-w)/2\) satisfies \(w\bar u=-u\). Replacing \(x\) by \(ux\)
normalizes the action to \(c(rx)=-\bar r x\).

For a basis \(e_j=\rho^jx\), a reflection \(\rho^bc\) sends
\(e_j\) to \(-e_{b-j}\). Its invariant lattice is generated by the
differences \(e_j-e_{b-j}\). Two different reflections give a connected
graph on the five indices: their composition is a nonzero translation.
Their differences therefore generate the augmentation ideal \(I_{\mathcal R}\).
Adding \(\Lambda^{C_5}=\Z_5N_Gx\), the quotient is
\(\mathcal R/(I_{\mathcal R}+\Z_5N_G)\simeq\Z/5\).
This computes the five-adic index of the actual subfield Bloch lattices.
Guo--Qin's bound \(Q_2\mid5\) makes it the full index.
\end{proof}

Apply Proposition \ref{prop:lattice} to the good lines to finish
Theorem \ref{thm:main}. Varying \(\ell\) changes the discriminant, so the
result is an infinite family rather than a sequence of repetitions.

\section{Fixed data and reproducibility}
The theorem uses a fixed finite extension to define its prime set.
It also has a small-field exact description. In \(E=\Q(y)\), \(y=\zeta_{20}\),
one may take the Tate-line generator
\[
 g_2=-5y^7+7y^6-3y^5+y^4-5y^3+6y^2-3.
\]
The supplied PARI/GP 2.15.5 certificate proves \(h(E)=1\), certifies the
fundamental unit group, and records \(g_2\)'s coordinates \((1,3,1,0)\).
Thus it is not a fifth power. If \(\sigma_{13}(y)=y^{13}\) and
\(c(y)=y^{-1}\), the unit coordinates of
\(\sigma_{13}(g_2)/g_2^2\) and \(c(g_2)/g_2\) are respectively
\((0,-5,-5,15)\) and \((0,0,0,10)\), where the final coordinate refers to
the order-twenty roots of unity. Both ratios are fifth powers. This gives the
required \(\Delta\)-character \(\omega^{-1}\). Since it is an \(S_0\) Kummer
class and the base cohomology is one-dimensional, it generates the canonical
line \(D\).

The fixed pure-wild field \(M_0\) is the splitting field of
\[
 f_0(X)=X^5+10X^3-10X^2-15X-18.
\]
A separate ray-class certificate over \(\Q(i)\) proves its group,
conductor, and discriminant. For \(\ell\equiv1\pmod{20}\), membership in
\(\mathcal P\) can be checked by
\[
 g_2^{(\ell-1)/5}\equiv1\pmod Q\quad(Q\mid\ell\text{ in }E),
 \qquad f_0\bmod\ell\text{ irreducible}.
\]
The second test detects the nontrivial cyclic Frobenius, since \(\ell\) splits
in the quadratic subfield. The exact certificate checks all eight residue
symbols in \(E\). It finds that \(701,941,1321,1381,1801,2081\) meet both
conditions, whereas \(41\) fails the Tate condition. These calculations check
fixed data and examples; the infinite-family proof uses Chebotarev and the
cup-product argument, not a finite search.

\section{Relation to earlier work and scope}
Browkin--Gangl \cite[Section 12.4]{BG} already gave numerical evidence for
\(Q_2=5\), including a dihedral Hilbert-class-field example with quadratic
subfield \(\Q(\sqrt{-47})\). Their numerical lattice-generation qualifications
must be kept distinct from an integral proof. Guo--Qin \cite{GQ} established
the dihedral divisibility theorem, which we use, and an infinite family with
\(Q_2=1\) and quadratic subfield \(\Q(\sqrt{-3})\).
General higher \(K\)-group index formulae and their cohomological interfaces
already occur in Caputo \cite{Caputo12,Caputo26} and Bartel--de Smit
\cite{BS}. The general integral
co-descent and genus statements used here belong to Assim--Movahhedi
\cite{AM12,AM19}; the generalized Bockstein theorem belongs to
Lam--Liu--Sharifi--Wake--Wang \cite{LLSWW}.

The arithmetic construction here keeps the quadratic subfield fixed and
selects among four mixed rays with a uniform two-root bound. It is not a new
general Bockstein theory or a new polynomial root-counting lemma. Nor does it
prove that every ray chosen by a separate Kummer algorithm is good. No claim
of a formally verified proof, a first numerical example, or a named open
problem's historical attribution is made.

\paragraph{Computational and drafting assistance.}
Language-model assistance was used in developing the draft and accompanying
exact-arithmetic scripts. The mathematical argument is supplied for independent
human scrutiny. The artifacts contain exact checks, not Lean formalization or
a certificate of novelty.

\begin{thebibliography}{99}
\bibitem{BG} J. Browkin and H. Gangl,
\emph{Tame kernels and second regulators of number fields and their subfields},
Journal of K-Theory \textbf{12} (2013), 137--165,
\href{https://doi.org/10.1017/is013005031jkt229}{doi:10.1017/is013005031jkt229}.
\bibitem{GQ} X. Guo and H. Qin,
\emph{The extended Bloch groups of biquadratic and dihedral number fields},
Journal of Pure and Applied Algebra \textbf{222} (2018), 3968--3981,
\href{https://doi.org/10.1016/j.jpaa.2018.02.015}{doi:10.1016/j.jpaa.2018.02.015}.
\bibitem{AM12} J. Assim and A. Movahhedi,
\emph{Norm index formula for the Tate kernels and applications},
Journal of K-Theory \textbf{9} (2012), 359--383,
\href{https://doi.org/10.1017/is011003006jkt135}{doi:10.1017/is011003006jkt135}.
\bibitem{AM19} J. Assim and A. Movahhedi,
\emph{Galois co-descent for motivic tame kernels},
Annales math\'ematiques du Qu\'ebec \textbf{49} (2025), 503--555,
\href{https://doi.org/10.1007/s40316-024-00233-8}{doi:10.1007/s40316-024-00233-8}.
The theorem numbers cited here refer to the
\href{https://arxiv.org/abs/1901.07219}{2019 preprint version}.
\bibitem{LLSWW} Y. H. J. Lam, Y. Liu, R. Sharifi, P. Wake, and J. Wang,
\emph{Generalized Bockstein maps and Massey products},
Forum of Mathematics, Sigma \textbf{11} (2023), e5,
\href{https://doi.org/10.1017/fms.2022.103}{doi:10.1017/fms.2022.103};
\href{https://arxiv.org/abs/2004.11510}{arXiv:2004.11510}.
\bibitem{Caputo12} L. Caputo,
\emph{The Brauer--Kuroda formula for higher S-class numbers in dihedral
extensions of number fields}, Acta Arithmetica \textbf{151} (2012), 217--239,
\href{https://doi.org/10.4064/aa151-3-1}{doi:10.4064/aa151-3-1}.
\bibitem{BS} A. Bartel and B. de Smit,
\emph{Index formulae for integral Galois modules},
Journal of the London Mathematical Society \textbf{88} (2013), 845--859,
\href{https://doi.org/10.1112/jlms/jdt033}{doi:10.1112/jlms/jdt033}.
\bibitem{Caputo26} L. Caputo, \emph{Regulator constants and cohomology},
Journal of the London Mathematical Society \textbf{114} (2026),
\href{https://doi.org/10.1112/jlms.70671}{doi:10.1112/jlms.70671};
\href{https://arxiv.org/abs/2603.01310}{arXiv:2603.01310}.
\bibitem{Levine} M. Levine,
\emph{The indecomposable \(K_3\) of fields},
Annales scientifiques de l'\'Ecole Normale Sup\'erieure (4)
\textbf{22} (1989), 255--344,
\href{https://doi.org/10.24033/asens.1585}{doi:10.24033/asens.1585}.
\end{thebibliography}
\end{document}
